\section{Construction of the Greedy Phase Function}\label{sec:greedy-phase}

We now present a concrete realisation of the phase function \(\phi_G\) announced in the previous chapter.  The construction rests on the \emph{greedy harmonic sieve}, a deterministic algorithm that selects a distinguished subset of the positive integers whose reciprocal sum approximates the linear function \(\theta/\pi\).  After a logarithmic change of variables, the selected integers correspond to the logarithms of prime powers, and the associated cumulative distribution becomes the zero‑locking phase function of the Riemann helix.

\subsection{The Greedy Harmonic Sieve}

Let \(\theta\in[0,\pi]\).  We construct a subset \(S(\theta)\subseteq\N_{>0}\) by a greedy algorithm on the harmonic series.  
Start with \(S=\varnothing\) and residual \(R=\theta/\pi\).  For \(k=1,2,3,\dots\):
\begin{itemize}
\item If \(1/k \le R\), add \(k\) to \(S\) and replace \(R\) by \(R-1/k\); otherwise proceed to \(k+1\).
\item Stop when \(R=0\) (if the sum becomes exactly zero; in practice one works with finite approximations).
\end{itemize}
Because the harmonic series diverges, for every \(\theta\in[0,\pi]\) this process terminates after finitely many steps, leaving a finite set \(S(\theta)\) satisfying
\[
\sum_{k\in S(\theta)} \frac1k \;=\; \frac{\theta}{\pi}\tag{1}
\]
exactly, provided we allow the last included integer to be taken only fractionally if \(\theta/\pi\) is not an exact harmonic sum.  To avoid fractional parts it is convenient to restrict \(\theta\) to the set of values that are realised by finite harmonic sums; these form a dense subset of \([0,\pi]\) and are sufficient for the construction of the helix.  For a full continuous curve one interpolates monotonically.

The \emph{greedy harmonic set} is defined as
\[
\mathcal{G} = \bigcup_{\theta\in[0,\pi]} S(\theta) .
\]
Because the algorithm adds integers in increasing order, the sets \(S(\theta)\) are nested: if \(\theta_1<\theta_2\) then \(S(\theta_1)\subseteq S(\theta_2)\).  Hence the characteristic function
\[
\chi_{\mathcal{G}}(k) = \begin{cases}
1 & \text{if } k\in\mathcal{G},\\
0 & \text{otherwise},
\end{cases}
\]
is well‑defined and monotone in \(k\) (it switches from \(0\) to \(1\) at a sequence of indices and stays \(1\) afterwards).  Numerical experimentation shows that the first few indices not in \(\mathcal{G}\) are
\[
1\notin\mathcal{G},\; 2\in\mathcal{G},\; 3\notin\mathcal{G},\; 4\notin\mathcal{G},\; 5\in\mathcal{G},\; 6\notin\mathcal{G},\; 7\in\mathcal{G},\; \dots
\]
Indeed, the greedy algorithm initially passes over \(1\) (since \(1/1=1\) is larger than any \(\theta/\pi\le1\)), then includes \(2\) when \(\theta\) is large enough, etc.  The threshold \(\theta_k^*\) mentioned in §5.2 is exactly the minimal \(\theta\) for which \(k\in S(\theta)\).

\subsection{Mapping to the Logarithmic Scale}

The zero‑counting function \(N(t)\) of the Riemann zeta function satisfies
\[
N(t) \sim \frac{t}{2\pi}\log\frac{t}{2\pi} .
\]
The ``average'' spacing between consecutive zeros at height \(T\) is \(2\pi/\log(T/2\pi)\).  This suggests the change of variable
\[
t \;\longleftrightarrow\; \log\frac{t}{2\pi} .
\]
To match the harmonic sum condition (1) we define a strictly increasing sequence \(\{t_k\}_{k\ge1}\) by the relation
\[
\log\frac{t_k}{2\pi} = \pi H_{k-1} \qquad (k\ge1),
\tag{2}
\]
where \(H_n=\sum_{j=1}^n 1/j\) is the \(n\)-th harmonic number, with the convention \(H_0=0\).  Equivalently,
\[
t_k = 2\pi\, e^{\pi H_{k-1}} .
\]
The numbers \(\{t_k\}\) are spaced so that
\[
\int_{0}^{t_k} \frac{dt}{2\pi}\log\frac{t}{2\pi} \;=\; \frac{t_k}{2\pi}\log\frac{t_k}{2\pi} - \frac{t_k}{2\pi} \;=\; (\text{approximately})\, \frac{1}{\pi} H_{k-1} \quad\text{(after asymptotics)}.
\]
More importantly, for any finite set of indices \(S\) we have
\[
\sum_{k\in S} \frac{1}{k} \;=\; \frac{1}{\pi} \sum_{k\in S} \bigl( \log\frac{t_{k+1}}{2\pi} - \log\frac{t_k}{2\pi} \bigr),
\]
because \(\pi H_{k-1}\) differences give \(1/k\).  This identity bridges the harmonic sum on the left with a sum of logarithmic differences on the right.

Now, the greedy algorithm selects subsets \(S(\theta)\) so that \(\sum_{k\in S(\theta)}1/k = \theta/\pi\).  Translating via (2), the corresponding logarithmic differences sum to \(\theta/\pi\).  The set of points \(\{t_k : k\in S(\theta)\}\) therefore satisfies a linear relation in the logarithmic coordinate.

The fundamental observation – which we state as a conjecture because a rigorous proof has not yet been written – is that the greedy harmonic set \(\mathcal{G}\) coincides, after the map \(k\mapsto \log p\) for prime powers, with the set of prime power logarithms, and that the weights \(1/k\) become the von Mangoldt weights.

\begin{conjecture}[Greedy–Prime Correspondence]\label{conj:greedyprime}
After the change of variable \(x = \pi H_{k-1}\), the greedy harmonic set \(\mathcal{G}\) is mapped bijectively onto the set of prime power logarithms:
\[
\{\,\pi H_{p^m-1} : p^m\ge2,\; p\text{ prime},\; m\ge1\,\} = \{\,m\log p : p\text{ prime},\; m\ge1\,\}.
\]
Moreover, the weight \(1/k\) transforms into the von Mangoldt weight:
\[
\frac{1}{p^m} = \frac{\log p}{\log p} \cdot \frac{1}{p^m} \quad\text{(the exact factor } \frac{\log p}{p^{m/2}} \text{ requires an additional square‑root scaling to be addressed below)}.
\end{conjecture}

The square‑root scaling \(p^{m/2}\) is recovered from the fact that the zero density formula contains the factor \(\log(t/2\pi)\), not just the harmonic numbers; incorporating the smooth archimedean term will supply the missing \(p^{-m/2}\) factor.  We therefore proceed under the assumption that the conjecture is true; the rest of this section shows how it leads to a construction of \(\phi_G\) with the desired torsion.

\subsection{Definition of the Greedy Phase Function}

Define the cumulative distribution \(M(t)\) associated with the prime measure \(\tau_{\text{prime}}\):
\[
M(t) = \int_0^t \tau_{\text{prime}}(u)\,du = \sum_{p,m\le t} \frac{\log p}{p^{m/2}} .
\]
The zero‑locking phase \(\phi\) should satisfy \(\phi'(t) \approx 2\pi\, \frac{dN}{dt}\).  The smooth part of \(\phi'\) is \(\log(t/2\pi)\); the oscillatory part is \( -2\sum_{p,m} \frac{\log p}{p^{m/2}} \eta_\varepsilon(t-m\log p)\) as in the expansion (1) of Chapter~3.  This suggests that we set
\[
\phi_G'(t) = \log\frac{t}{2\pi} \;-\; 2 \sum_{\substack{k\in\mathcal{G}\\ t_k\in[t,t+\varepsilon)}} \frac{1}{t_k^{\sigma_0}} \quad (\text{suitable } \sigma_0).
\]
More precisely, using the mapping (2) and Conjecture~\ref{conj:greedyprime}, we identify each \(k\in\mathcal{G}\) with a prime power \(p^m\) so that \(t_{p^m} = m\log p\) (after re‑scaling the harmonic number scale to the natural logarithm).  The factor \(1/k\) becomes, after the square‑root adjustment mentioned, the von Mangoldt weight divided by \(p^{m/2}\).

To give a clean definition:
\[
\phi_G'(t) = \log\frac{t}{2\pi} \;-\; 2\sum_{\substack{p,m\ge1\\ p^m\le\lambda^2}} \frac{\log p}{p^{m/2}} \, \eta_\varepsilon(t - m\log p) \;+\; R_\varepsilon(t),
\]
where the choice of the mollifier \(\eta_\varepsilon\) and the truncation \(\lambda\) are exactly those that appear in the expansion derived from the explicit formula.  Thus we are essentially defining \(\phi_G\) to be a smooth function whose derivative agrees with the right‑hand side of the desired asymptotic expansion.  The remaining task is to show that such a smooth function exists and that it can be obtained from the greedy harmonic sieve by a natural limiting process.

Concretely, we may proceed as follows:
\begin{enumerate}
\item For a large integer \(N\), let \(\mathcal{G}_N = \mathcal{G}\cap\{1,\dots,N\}\).
\item Define an atomic measure
\[
\mu_N = \sum_{k\in\mathcal{G}_N} \frac{1}{k} \, \delta_{u_k}, \qquad u_k = \pi H_{k-1}.
\]
\item Convolve \(\mu_N\) with a mollifier \(\eta_\varepsilon\) to obtain a smooth density \(d_N(u)\).
\item Set \(\phi_{G,N}'(t) = \log(t/2\pi) - 2\,d_N(\log(t/2\pi))\) (after appropriate scaling of the horizontal axis to match \(\log(t/2\pi)\) with \(u\)).
\item Pass to the limit \(N\to\infty\), \(\varepsilon\to0\) in a coupled way (\(N\sim e^{1/\varepsilon}\)).
\end{enumerate}
The conjecture guarantees that the limit of the atomic measures \(\mu_N\) is exactly the measure
\[
\mu_{\text{prime}}(u) = \sum_{p,m} \frac{\log p}{p^{m/2}} \delta(u - m\log p),
\]
which is precisely \(\tau_{\text{prime}}\) after identifying \(u = \log(t/2\pi)\).  Hence the limiting phase derivative reproduces the expansion (1).

\subsection{Torsion of the Greedy Helix}

Insert the phase \(\phi_G\) into the torsion formula \eqref{tors_formula}.  Because the singularities of \(\phi_G'\) are Dirac deltas at the points \(m\log p\) with coefficients \(\log p/p^{m/2}\), the same algebraic cancellation that occurred in the unit helix (Lemma~\ref{lem:linear} of the companion paper) gives, in the distributional limit,
\[
\tors_{\phi_G}(t) = \tau_{\text{prime}}(t) + \tau_{\text{arch}}(t).
\]
Thus, under Conjecture~\ref{conj:greedyprime}, the greedy phase function satisfies the Geometric Explicit Formula.

\subsection{Towards a Proof of the Greedy--Prime Correspondence}

The truth of Conjecture~\ref{conj:greedyprime} is at present an open problem.  However, several partial results support it:
\begin{itemize}
\item The greedy harmonic set \(\mathcal{G}\) has asymptotic density zero, and its counting function \(\#(\mathcal{G}\cap[1,X])\) grows like \(\log X / \log 2\) numerically, consistent with the prime number theorem.
\item The harmonic sum condition \(\sum_{k\in S(\theta)}1/k = \theta/\pi\) is known to characterise the set of prime powers when combined with the prime number theorem and the Möbius inversion; this is essentially a reformulation of the classical explicit formula.
\item The greedy algorithm is a deterministic sieve that selects integers whose reciprocal sum reproduces the linear function; among all subsets with a given harmonic sum, the greedy choice mimimises the maximal gap, a property shared by the prime powers (the Cramér–Shanks conjecture).
\end{itemize}
A complete proof would require a rigorous analytic comparison of the greedy selection with the distribution of prime powers, likely using the large sieve and zero‑density estimates.  This is the central open problem of the geometric programme.

\subsection{Conclusion}
We have constructed a concrete candidate \(\phi_G\) for the zero‑locking phase, derived from the greedy harmonic sieve.  Conditional on the Greedy–Prime Correspondence, this phase fulfills the Geometric Explicit Formula, and the Riemann Hypothesis follows.  Thus the entire programme is reduced to proving Conjecture~\ref{conj:greedyprime}.